A frozen SciGeoGuard-X verifier was evaluated on 445 controlled
SilentGeoBench workflows (145 PROVEN,
278 REFUTED, and 22 UNPROVEN).
These generated reference labels test consistency with encoded scientific invariants and are
not independent human ground truth. Against the same generated reference, the decomposed
Qwen2.5-7B baseline achieved 44.44\% accuracy with an FSCR of 69.64\%,
while Mistral-7B achieved 38.33\% accuracy with an FSCR of 69.64\%.

Live-data counterfactuals were then evaluated in Bengaluru, Rotterdam, California's Central
Valley, and Manaus. Replacing Sentinel-2 B08 NIR with B11 SWIR in an NDVI-labelled
calculation produced a regional median mean absolute index difference of
0.0977 (95\% region-bootstrap CI 0.0732--0.1242). Bilinear interpolation of
ESA WorldCover nominal classes produced a median
12.99\% invalid class values
across regions (95\% CI 7.49\%--19.06\%), whereas nearest-neighbour resampling preserved valid class codes in
100\% of evaluated outputs.
For Copernicus DEM GLO-30, threefold interpolation along each raster axis generated
9 times as many grid cells, so
88.89\% of fine-grid positions did not correspond
to distinct source samples; the corrected region-median core p95 native-to-fine-to-native
round-trip error was 1.3998 m.
Finally, live rainy GPM IMERG V07 cases were recovered in all
4 study regions. Omitting the 0.5 h duration when integrating
half-hourly precipitation rates yielded the expected
2.0$\times$ numeric overstatement, while treating the
full-day mean rate as daily accumulation yielded
0.041667 (=1/24) of the correct accumulation.